\section{관련 연구}
\label{sec:관련연구}

본 장에서는 FANET 환경을 위해 제안된 기존의 라우팅 기법과 기계학습 및 정책 증류의 관련 연구 동향을 요약하고 제안 기법과의 차별성을 밝힌다.

\subsection{지리적 및 예측적 라우팅 프로토콜}
무인 항공기 기반 네트워크의 높은 가동성을 고려한 지리적 라우팅 기법 중 GPSR(Greedy Perimeter Stateless Routing)\cite{gpsr}은 이웃 노드 중 목적지에 물리적 거리가 가장 가까운 노드를 차홉으로 선별하는 대표적인 기법이다. 그러나 지리적 라우팅 프로토콜은 노드가 배치되지 않은 장애물이나 통신 불가능 영역인 '라우팅 홀(Routing Hole)'에 가로막힐 경우 로컬 미니멈(Local Minimum) 상태에 빠지게 된다. 이를 극복하기 위해 패킷의 이동 경로 평면화(Planarization) 과정을 거쳐 경계면을 따라 우회 탐색하는 우회 모드(Perimeter Mode)를 사용하나, 노드의 고속 이동성으로 인해 경계면 토폴로지가 컴파일 타임 이전에 파괴되는 현상으로 인해 패킷 손실이 매우 흔하게 발생한다. 이를 보완하기 위해 칼만 필터나 이동 방향의 기하학적 보정을 통해 미래 노드 위치를 예측하고 포워딩 결정을 돕는 예측적 지리적 라우팅 기법들이 고안되었으나, 무선 무작위 채널의 비정상성이나 정교하게 고안된 토폴로지 장애물 상황에서 최적의 파라미터를 찾는 데 명확한 한계를 지닌다.

\subsection{강화학습 및 다중 에이전트 강화학습 기반 라우팅}
최근 무선 애드혹 환경의 동적 링크 상태 변화를 Markov 의사결정 과정으로 간주하고 이를 학습 기반으로 풀고자 하는 강화학습 라우팅 기법이 급격히 증가하고 있다. 대표적으로 단일 에이전트 Q-learning 모델에 멀티홉 Q값 결정을 결합한 IQMR Q($\lambda$)\cite{iqmr}이 제안되었으며, 링크의 시계열 진화 경향성을 Q학습과 결합하여 우회 경로의 일반화를 도모한 Evo-QGeo\cite{evoqgeo} 등이 제시되었다. 
나아가, 여러 대의 UAV가 스스로 무선 링크 형성 전략을 공조 학습하는 다중 에이전트 강화학습(MARL) 모델도 활발히 연구되고 있다. 일례로 다중 에이전트 간의 학습된 임시 메시지(Emergent Communication) 교환을 통해 경로 혼잡도와 대안 라우팅을 협조적으로 탐색하는 DRAMA\cite{drama}와, LoRa 기반 FANET 환경에서 에너지 효율적 노드 군집을 위해 GLo-MAPPO\cite{glo_mappo}를 접목하는 연구 등이 대표적이다. 
그러나 이들 MARL 및 학습 기반 라우팅 기법은 각 UAV가 매번 이웃 에이전트들과 대규모의 토폴로지 및 신경망 벡터 정보(Control plane bytes)를 무선 교환해야만 차홉 추론이 가능하므로, 제어 채널의 과부하와 간섭을 동반하여 실질 배포가 불가능한 배포 격차 문제를 안고 있다.

\subsection{정책 증류 및 지식 전수 기법}
정책 증류(Policy Distillation)\cite{rusu2016policy}는 강화학습을 통해 성공적으로 훈련된 교사(Teacher) 정책의 지식을 가벼운 학생(Student) 네트워크로 이전하는 학습 프레임워크로, 주로 강화학습의 연산 비용 압축이나 지식 전수를 위해 사용된다. 이공계 시스템 전반에서 CTDE(Centralized Training, Decentralized Execution) 철학을 극대화하기 위해, 학습 단계에서는 전역 상태 정보를 온전히 관측하는 Privileged Teacher를 학습시킨 뒤, 실행 단계에서는 오직 에이전트의 로컬 부분관측 정보만으로 구동되는 Student 모델로 정책을 증류하는 기법이 주목받고 있다. 
본 논문은 이러한 정책 증류 기법을 FANET 라우팅 문제에 최초로 대입하여, 전역 그래프 상태 정보를 모두 파싱하여 최적의 지름길을 아는 GNN 기반 교사의 능력을, 1-hop 관측 정보만 가볍게 수집하여 실행 오버헤드가 거의 없는 MLP 기반 학생으로 오프라인 전수한다. 아울러 예측적 위험 스위칭이라는 동적 트리거링 구조와 결합하여 학습 모델 자체의 연산 부담을 추가적으로 제어한다는 점에서 기존 연구들과 명확한 독창적 차별점을 지닌다.
